\documentclass[12pt]{article}

% Packages
\usepackage[margin=2.5cm]{geometry}
\usepackage{amsmath}
\usepackage{amssymb}
\usepackage{bm}
\usepackage{enumitem}
\usepackage{graphicx}
\usepackage[hidelinks]{hyperref}

% No indentation
\setlength{\parindent}{0pt}
\setlength{\parskip}{6pt}

\begin{document}

\title{A Practical Model-Agnostic Protocol for Extracting\\
$Q_{\max}$ and $K_{\mathrm{aff}}$ from Adsorption Data\\[6pt]
\large (Version 3 --- with validation figures)}
\author{}
\date{}
\maketitle

% ======================================================================
\section{Scope and minimal assumptions}
% ======================================================================

We consider equilibrium adsorption from aqueous solution at fixed
temperature $T$ and fixed chemistry (pH, ionic strength, background ions,
cosolutes).  Let $C\ge 0$ be the equilibrium aqueous concentration and
$q(C)$ the equilibrium uptake per adsorbent mass.

A model-agnostic extraction of $Q_{\max}$ and an effective affinity
constant $K_{\mathrm{aff}}$ is possible under the following assumptions:

\begin{itemize}
  \item[A1.] \textbf{Independent-site superposition.}  The adsorption
        can be represented as a superposition of independent Langmuir
        site classes (heterogeneous energies allowed, no lateral
        interactions or cooperativity in the site-filling step).
        This is the standard integral equation of adsorption on
        heterogeneous surfaces~\cite{Sips1948,Jaroniec1988,Rudzinski1992}.
  \item[A2.] \textbf{Finite saturating capacity.}  A finite saturating
        capacity exists for the saturating part of adsorption, i.e.,
        $Q_{\max}=\int_0^\infty\rho(K)\,dK<\infty$.
  \item[A3.] \textbf{Finite first moment (Henry constant).}  The first
        moment $\int_0^\infty K\rho(K)\,dK<\infty$ is finite, so that a
        well-defined Henry constant exists.
        \emph{Note:} This condition is violated by pure Freundlich
        systems ($f(K)\propto K^{n-1}$, $n<1$) and Sips systems with
        $m<1$ in the thermodynamic limit.  In practice, any real
        adsorbent has a finite maximum affinity, so $K_H$ is always
        finite experimentally, but it may be very sensitive to the
        high-affinity tail.  See Section~\ref{sec:convergence}.
  \item[A4.] \textbf{Sufficient experimental coverage.}  The dataset
        contains points in the Henry regime ($C\to 0$) and in a
        high-concentration regime that either approaches a plateau or
        can be extrapolated to a plateau.
\end{itemize}

\textbf{Scope of A1.}  Assumption A1 excludes cooperative binding
(Hill isotherm with $n_H>1$), adsorbate--adsorbate lateral interactions
(Fowler--Guggenheim), and multilayer adsorption (BET).  The protocol
should \emph{not} be applied to sigmoidal isotherms or isotherms with
inflection points, as these indicate mechanisms outside the heterogeneous
Langmuir class.

For PFAS on many solids, a second contribution is sometimes present: a
linear partition-like term that does not saturate within the measured
range~\cite{Vieth1976}.  We therefore treat two cases:
(i)~purely saturating adsorption, and
(ii)~partition plus saturating adsorption.

% ======================================================================
\section{General heterogeneous Langmuir representation}
% ======================================================================

\subsection{Purely saturating heterogeneous adsorption}

We represent heterogeneity by a nonnegative ``capacity density''
$\rho(K)$ over affinity constants $K>0$~\cite{Sips1948,Jaroniec1988}:
\begin{equation}
  q(C) = \int_{0}^{\infty} \rho(K)\,\frac{KC}{1+KC}\,dK,
  \qquad \rho(K)\ge 0.
  \label{eq:het_langmuir}
\end{equation}
Here $\rho(K)\,dK$ is the saturation capacity carried by site classes
with affinities in $[K,K+dK]$.  The normalised distribution is
\begin{equation}
  Q_{\max} \equiv \int_0^\infty \rho(K)\,dK,
  \qquad
  f(K) \equiv \frac{\rho(K)}{Q_{\max}},
  \qquad
  \int_0^\infty f(K)\,dK = 1.
  \label{eq:f_def}
\end{equation}

\subsection{Partition plus saturating adsorption}

If a linear, nonsaturating contribution is present (partition into
organic matter, swelling domains, or other linear
uptake)~\cite{Vieth1976,Paul1976}, we write
\begin{equation}
  q(C) = K_p\,C + \int_{0}^{\infty} \rho(K)\,\frac{KC}{1+KC}\,dK,
  \qquad K_p \ge 0.
  \label{eq:partition_plus_langmuir}
\end{equation}

% ======================================================================
\section{Invariants: $Q_{\max}$, $K_H$, and $K_{\mathrm{aff}}$}
% ======================================================================

\subsection{Saturation capacity}

For the purely saturating case, since $KC/(1+KC)\to 1$ as
$C\to\infty$ for every fixed $K>0$, dominated convergence gives
\begin{equation}
  \lim_{C\to\infty} q(C)
  = \int_0^\infty \rho(K)\,dK = Q_{\max}.
  \label{eq:Qmax_limit}
\end{equation}

\subsection{Henry constant under heterogeneity}

For small $C$, the Langmuir kernel satisfies
$KC/(1+KC) = KC + O(C^2)$.
If Assumption~A3 holds,
differentiation under the integral gives~\cite{Jaroniec1988,Rudzinski1992}
\begin{equation}
  K_H \equiv \left.\frac{dq}{dC}\right|_{C\to 0}
  = \int_0^\infty K\,\rho(K)\,dK
  = Q_{\max}\int_0^\infty K f(K)\,dK.
  \label{eq:KH_moment}
\end{equation}

In the partition-plus-saturating case,
$dq/dC|_{C\to 0} = K_p + K_H$,
so one must separate $K_p$ if present.

\subsection{Effective affinity constant}

A single effective affinity constant consistent with both the saturating
capacity and the Henry slope is~\cite{Klotz1971,Nederlof1990}
\begin{equation}
  K_{\mathrm{aff}} \equiv \frac{K_H}{Q_{\max}}
  = \int_0^\infty K\,f(K)\,dK,
  \label{eq:Kaff_def}
\end{equation}
i.e., $K_{\mathrm{aff}}$ is the capacity-weighted mean affinity over the
(unknown) site distribution.

\subsection{Convergence conditions and limitations}
\label{sec:convergence}

\textbf{When $K_H$ diverges.}
For the Freundlich isotherm $q=K_F C^{1/n}$ ($0<1/n<1$), the underlying
distribution $f(K)\propto K^{1/n-1}$ is a power law whose first moment
diverges~\cite{Jaroniec1988}.
Operationally, one can define
\begin{equation}
  K_H^{\mathrm{oper}} \equiv
  \frac{\sum_{i\in\mathcal{I}_0} w_i C_i q_i}
       {\sum_{i\in\mathcal{I}_0} w_i C_i^2},
  \label{eq:KH_operational}
\end{equation}
but this ``Henry constant'' depends on $\mathcal{I}_0$ and is not a
material constant in the strict thermodynamic sense.

\textbf{When $A_1$ diverges.}
The high-$C$ expansion (Eq.~\ref{eq:highC_expansion}) requires
$A_1=\int_0^\infty\rho(K)/K\,dK<\infty$.  If $\rho(K)$ has substantial
weight near $K=0$, the $1/C$ extrapolation converges slowly.

% ======================================================================
\section{Practical workflow}
\label{sec:workflow}
% ======================================================================

This section describes a robust workflow that does not assume a specific
$f(K)$, and returns $Q_{\max}$ and $K_{\mathrm{aff}}$ (and optionally
$K_p$).

\subsection{Step 0: Data preparation}

Let the dataset be $\{(C_i,q_i,\sigma_i)\}_{i=1}^N$, sorted by
$C_i$ in ascending order.  Define weights $w_i\equiv 1/\sigma_i^2$.

\subsection{Step 1: Decide whether a partition term is needed}

\begin{enumerate}[nosep]
  \item Select the highest-$C$ subset $\mathcal{I}_\infty$ (last $m\ge 4$ points).
  \item Fit $q_i = \beta_0 + \beta_1 C_i$ by WLS over $\mathcal{I}_\infty$.
  \item Compute $t = \widehat{\beta}_1 / \mathrm{SE}(\widehat{\beta}_1)$
        and test $H_0:\beta_1=0$ at $\alpha=0.05$.
  \item If rejected and $\widehat{\beta}_1>0$: set $K_p>0$.
  \item Otherwise: set $K_p=0$.
\end{enumerate}

\subsection{Step 2: Estimate $Q_{\max}$}

\subsubsection{Case A: Plateau observed}

Weighted average over the plateau region $\mathcal{I}_\infty$:
\begin{equation}
  \widehat{Q}_{\max} =
  \frac{\sum_{i\in\mathcal{I}_\infty} w_i q_i}
       {\sum_{i\in\mathcal{I}_\infty} w_i}.
  \label{eq:Qmax_plateau}
\end{equation}

\textbf{Selection of $\mathcal{I}_\infty$:}
Start from the last $m_0=3$ points and increase $m$ until the slope of
$q$ vs.\ $\ln C$ becomes significant at $\alpha=0.05$.

\subsubsection{Case B: $1/C$ extrapolation}

\begin{equation}
  q(C) = Q_{\max} - \frac{A_1}{C} + O\!\left(\frac{1}{C^2}\right),
  \qquad
  A_1 \equiv \int_0^\infty \frac{\rho(K)}{K}\,dK.
  \label{eq:highC_expansion}
\end{equation}

Fit $q_i \approx \beta_0 + \beta_1/C_i$ and set
$\widehat{Q}_{\max}=\widehat{\beta}_0$.
Diagnostic: check $\widehat{A}_1=-\widehat{\beta}_1 > 0$.

\subsection{Step 3: Estimate $K_H$ (and $K_p$ if needed)}

Estimate $K_H$ by WLS over $\mathcal{I}_0$:
\begin{equation}
  \widehat{K}_H =
  \frac{\sum_{i\in\mathcal{I}_0} w_i C_i q_i}
       {\sum_{i\in\mathcal{I}_0} w_i C_i^2}.
  \label{eq:KH_fit}
\end{equation}

\textbf{Selection of $\mathcal{I}_0$ ($F$-test for curvature):}
\begin{enumerate}[nosep]
  \item Start with $\mathcal{I}_0=\{1,2,3\}$.
  \item For each candidate, fit $q=K_H C$ and $q=K_H C + \alpha C^2$.
  \item Apply an $F$-test; set $\mathcal{I}_0$ as the largest subset for
        which curvature is not significant at $\alpha=0.05$.
\end{enumerate}

If partition is present, estimate $K_p$ by WLS over $\mathcal{I}_\infty$
and iterate until $|\Delta K_p/K_p|<10^{-4}$ or 20~iterations.

\subsection{Step 4: Compute $K_{\mathrm{aff}}$}

\begin{equation}
  \widehat{K}_{\mathrm{aff}} = \frac{\widehat{K}_H}{\widehat{Q}_{\max}}.
  \label{eq:Kaff_est}
\end{equation}

\subsection{Step 5: Bootstrap uncertainty}

Resample $B=2000$ times with replacement, repeat Steps~1--4, and report
2.5th--97.5th percentile confidence intervals.

% ======================================================================
\section{Validation on synthetic data}
\label{sec:validation}
% ======================================================================

To validate the protocol, we generate synthetic data from four
representative models, add 5\% relative Gaussian noise, and apply the
workflow.  The models are:

\begin{enumerate}[nosep]
  \item \textbf{Langmuir:}
        $q = Q_{\max} K C/(1+KC)$, with $Q_{\max}=100$, $K=0.5$.
  \item \textbf{Sips (Langmuir--Freundlich):}
        $q = Q_{\max}(KC)^m/(1+(KC)^m)$, with $Q_{\max}=80$, $K=0.3$, $m=0.7$.
  \item \textbf{Partition + Langmuir:}
        $q = K_p C + Q_{\max} K C/(1+KC)$, with $Q_{\max}=50$, $K=1.0$, $K_p=2.0$.
  \item \textbf{Double Langmuir:}
        $q = Q_1 K_1 C/(1+K_1 C) + Q_2 K_2 C/(1+K_2 C)$, with
        $Q_1=60$, $K_1=5.0$, $Q_2=40$, $K_2=0.05$.
\end{enumerate}

All figures are generated by the MATLAB script
\texttt{universal\_route\_corrected.m}.

% ------ Figure 1 ------
\subsection{Isotherm reconstruction and region identification}

Figure~\ref{fig:isotherms} shows the four synthetic datasets together
with (i)~the true generating model (grey line), (ii)~the noisy data
(black circles), (iii)~the reconstructed effective Langmuir isotherm
using the estimated $\widehat{Q}_{\max}$ and $\widehat{K}_{\mathrm{aff}}$
(red dashed line), (iv)~the estimated $\widehat{Q}_{\max}$ (green
horizontal line), and (v)~the Henry slope $K_H C$ (blue dotted line).
The blue-shaded region marks the automatically selected Henry subset
$\mathcal{I}_0$, while the orange-shaded region marks the plateau
subset $\mathcal{I}_\infty$.

For the Langmuir model (panel~a), the reconstruction nearly overlaps
the true curve because the effective single-site Langmuir is exact.
For the Sips model (panel~b), the effective Langmuir captures the
correct $Q_{\max}$ and low-$C$ slope but deviates in the intermediate
range---this is expected because the Sips isotherm involves
heterogeneity that a single effective Langmuir cannot represent.
The partition-plus-Langmuir case (panel~c) correctly identifies $K_p>0$
and separates the saturating component.
The double Langmuir (panel~d) shows that $K_{\mathrm{aff}}$ is a
weighted average of the two true affinities, and $Q_{\max}$ correctly
recovers the sum $Q_1+Q_2$.

\begin{figure}[htbp]
  \centering
  \includegraphics[width=\textwidth]{fig_isotherms.pdf}
  \caption{Isotherm reconstruction for four synthetic models.
  Blue shading: Henry region $\mathcal{I}_0$ (automatically selected by
  $F$-test).  Orange shading: plateau region $\mathcal{I}_\infty$
  (selected by $t$-test on slope vs.\ $\ln C$).  Grey: true model;
  black circles: noisy data; red dashed: reconstructed effective
  Langmuir; green: $\widehat{Q}_{\max}$; blue dotted: Henry slope.}
  \label{fig:isotherms}
\end{figure}

% ------ Figure 2 ------
\subsection{Nonparametric affinity distribution recovery}

Figure~\ref{fig:nnls} shows the affinity distributions $\rho(K)$
recovered by regularised NNLS (Tikhonov with second-difference
smoothness, $\lambda$ selected by L-curve).  Red dashed vertical lines
mark the true affinity constants; the black dotted line marks the
computed $K_{\mathrm{aff}}$.

For the Langmuir model (panel~a), the recovered distribution is
concentrated around the true $K=0.5$, as expected for a homogeneous
surface.
For the Sips model (panel~b), the distribution is broader, reflecting
the heterogeneity encoded in $m=0.7<1$.  No single true $K$ exists in
this case.
The partition-plus-Langmuir case (panel~c) recovers a peak near the
true $K=1.0$ after subtracting the linear partition contribution.
The double Langmuir (panel~d) successfully resolves the two distinct
site populations at $K_1=5.0$ and $K_2=0.05$, separated by two
orders of magnitude.  The computed $K_{\mathrm{aff}}$ (dotted line)
lies between the two peaks, as it represents the capacity-weighted
mean affinity.

\begin{figure}[htbp]
  \centering
  \includegraphics[width=\textwidth]{fig_nnls_distributions.pdf}
  \caption{Nonparametric affinity distributions recovered by
  Tikhonov-regularised NNLS.  Blue bars: recovered $q_{m,j}$ on a
  $\log_{10}K$ grid; red dashed: true affinity constant(s); black
  dotted: computed $K_{\mathrm{aff}}$.}
  \label{fig:nnls}
\end{figure}

% ------ Figure 3 ------
\subsection{Sensitivity of $K_{\mathrm{aff}}$ to the Henry region}

Figure~\ref{fig:sensitivity} examines the robustness of
$\widehat{K}_{\mathrm{aff}}$ to the choice of the Henry region size
$|\mathcal{I}_0|$.  For each model, the number of low-$C$ points
included in $\mathcal{I}_0$ is varied from 3 to 15 and the resulting
$\widehat{K}_{\mathrm{aff}}$ is plotted.  Vertical dotted lines mark
the automatically selected $n_{\mathrm{henry}}$ (from the $F$-test).

For the Langmuir and double Langmuir models, $K_{\mathrm{aff}}$ is
nearly constant across all choices of $\mathcal{I}_0$, confirming that
the first moment of $\rho(K)$ is well-defined and the Henry slope is
robust.

For the Sips model, $K_{\mathrm{aff}}$ shows a systematic decrease as
more points are included.  This is the signature of a divergent first
moment: the Sips distribution with $m=0.7<1$ has a power-law tail, so
the operational $K_H$ depends on the concentration window.  This
behaviour serves as a diagnostic: if $K_{\mathrm{aff}}$ varies
substantially with $|\mathcal{I}_0|$, the underlying distribution
likely has a heavy tail and $K_{\mathrm{aff}}$ should be reported
together with the concentration range used.

\begin{figure}[htbp]
  \centering
  \includegraphics[width=0.75\textwidth]{fig_sensitivity.pdf}
  \caption{Sensitivity of $\widehat{K}_{\mathrm{aff}}$ to the number
  of points in the Henry region $\mathcal{I}_0$.  Vertical dotted
  lines: automatically selected $n_{\mathrm{henry}}$.  A flat curve
  indicates a well-defined first moment; a decreasing trend suggests
  a heavy-tailed distribution where $K_H$ is not a strict material
  constant.}
  \label{fig:sensitivity}
\end{figure}

% ------ Figure 4 ------
\subsection{Bootstrap uncertainty quantification}

Figure~\ref{fig:bootstrap} presents histograms of
$\widehat{Q}_{\max}$ (top row) and $\widehat{K}_{\mathrm{aff}}$
(bottom row) from 1000 bootstrap resamples.  Red lines indicate the
true values; black dashed lines indicate the point estimates from the
original dataset.

For all four models, the true values fall within the bootstrap
distributions, confirming that the protocol is unbiased.  The Langmuir
model (panel~a) shows the narrowest confidence intervals, as expected
for a homogeneous system with a well-defined single affinity.  The Sips
and double Langmuir models show broader distributions for
$K_{\mathrm{aff}}$, reflecting the additional uncertainty introduced by
heterogeneity.

The partition-plus-Langmuir model (panel~c) demonstrates that the
iterative $K_p$ estimation does not introduce systematic bias:
$Q_{\max}$ and $K_{\mathrm{aff}}$ are both well centred on their true
values.

\begin{figure}[htbp]
  \centering
  \includegraphics[width=\textwidth]{fig_bootstrap.pdf}
  \caption{Bootstrap distributions ($B=1000$).  Top row:
  $\widehat{Q}_{\max}$; bottom row: $\widehat{K}_{\mathrm{aff}}$.
  Red vertical lines: true values; black dashed: point estimates.}
  \label{fig:bootstrap}
\end{figure}

% ------ Figure 5 ------
\subsection{L-curve for regularization parameter selection}

Figure~\ref{fig:lcurve} illustrates the L-curve method for selecting
the Tikhonov regularization parameter $\lambda$ in the NNLS
distribution recovery (shown here for the Langmuir model).

The horizontal axis measures the residual norm
$\|\mathbf{r}\|=\|W^{1/2}(\mathbf{q}-A\mathbf{q}_m)\|$, which
decreases as $\lambda$ decreases (less regularization, better data fit).
The vertical axis measures the solution roughness
$\|\mathbf{L}\mathbf{q}_m\|$, which decreases as $\lambda$ increases
(more regularization, smoother distribution).

The optimal $\lambda^*$ is located at the corner of maximum curvature
(red circle), representing the best trade-off between fitting the data
and recovering a smooth, physically plausible distribution.  Too little
regularization (lower-left) produces a noisy, spiky distribution that
overfits the data; too much regularization (upper-right) produces an
overly smooth distribution that misses genuine features.

\begin{figure}[htbp]
  \centering
  \includegraphics[width=0.75\textwidth]{fig_lcurve.pdf}
  \caption{L-curve for the Langmuir model.  Each point corresponds to a
  different $\lambda$.  The red circle marks $\lambda^*$ at the corner of
  maximum curvature, selected automatically by the algorithm.}
  \label{fig:lcurve}
\end{figure}

% ======================================================================
\section{Optional: Nonparametric distribution recovery}
% ======================================================================

If you want to fit the entire curve without assuming a parametric $f(K)$,
discretize the integral on a logarithmic
grid~\cite{Umpleby2001,Umpleby2004,Nederlof1990,Cernik1995}.

Choose $M$ affinities on a log-spaced grid:
\[
  K_{\min} = \frac{0.1}{C_{\max}},\qquad
  K_{\max} = \frac{10}{C_{\min}},
\]
and approximate Eq.~\eqref{eq:het_langmuir} by
\begin{equation}
  q(C_i) \approx K_p C_i + \sum_{j=1}^{M} q_{m,j}\,\frac{K_j C_i}{1+K_j C_i},
  \qquad q_{m,j}\ge 0.
  \label{eq:disc_mix}
\end{equation}

Solve the Tikhonov-regularised NNLS
problem~\cite{Tikhonov1977,Hansen1998}:
\begin{equation}
  \min_{\bm{q}_m\ge 0,\,K_p\ge 0}\;
  \sum_{i=1}^{N} w_i\left(q_i - K_p C_i
    - \sum_{j=1}^{M} A_{ij}\,q_{m,j}\right)^2
  + \lambda \|\bm{L}\bm{q}_m\|_2^2,
  \label{eq:nnls_reg}
\end{equation}
where $A_{ij}=K_j C_i/(1+K_j C_i)$ and $\bm{L}$ is a second-difference
matrix in $\log K$ space.

\subsection{Regularization parameter selection}

The parameter $\lambda$ controls data fidelity vs.\ smoothness.
Standard approaches~\cite{Hansen1992}:
\begin{enumerate}[nosep]
  \item \textbf{L-curve method:} Plot $\log\|\bm{r}\|$ vs.\
        $\log\|\bm{L}\bm{q}_m\|$ and choose $\lambda$ at the corner
        of maximum curvature (Figure~\ref{fig:lcurve}).
  \item \textbf{GCV:} Minimize
        $V(\lambda)=\|W^{1/2}\bm{r}\|^2/(N-\mathrm{tr}(H_\lambda))^2$.
  \item \textbf{Discrepancy principle:} Choose $\lambda$ such that
        $\|\bm{r}\|^2\approx N\sigma^2$.
\end{enumerate}

\subsection{Consistency check}

After solving, recover
\begin{equation}
  \widehat{Q}_{\max}^{\mathrm{NNLS}} = \sum_{j=1}^{M}\widehat{q}_{m,j},
  \qquad
  \widehat{K}_H^{\mathrm{NNLS}} = \sum_{j=1}^{M} K_j\widehat{q}_{m,j},
  \qquad
  \widehat{K}_{\mathrm{aff}}^{\mathrm{NNLS}}
    = \frac{\widehat{K}_H^{\mathrm{NNLS}}}{\widehat{Q}_{\max}^{\mathrm{NNLS}}}.
\end{equation}
Compare with the direct asymptotic estimates from Steps~2--4.
Agreement within bootstrap confidence intervals provides a valuable
consistency check.

\subsection{Why the asymptotic method is the recommended default}

The nonparametric reconstruction is an ill-conditioned inverse problem
(Fredholm integral equation of the first kind), so the recovered
distribution depends on the regularisation~\cite{Hansen1998}.  In
contrast, $Q_{\max}$ and $K_H$ (hence $K_{\mathrm{aff}}$) are
directly anchored by the $C\to\infty$ and $C\to 0$ limits and are
substantially more robust.

% ======================================================================
\section*{References}
% ======================================================================

\begin{thebibliography}{99}

\bibitem{Sips1948}
R.~Sips, ``On the structure of a catalyst surface,''
\textit{J.\ Chem.\ Phys.}\ \textbf{16}, 490--495 (1948).

\bibitem{Jaroniec1988}
M.~Jaroniec and R.~Madey,
\textit{Physical Adsorption on Heterogeneous Solids},
Elsevier, Amsterdam (1988).

\bibitem{Rudzinski1992}
W.~Rudzinski and D.H.~Everett,
\textit{Adsorption of Gases on Heterogeneous Surfaces},
Academic Press, London (1992).

\bibitem{Klotz1971}
I.M.~Klotz and D.L.~Hunston,
``Properties of graphical representations of multiple classes of
binding sites,''
\textit{Biochemistry} \textbf{10}, 3065--3069 (1971).

\bibitem{Nederlof1990}
M.M.~Nederlof, W.H.~van Riemsdijk, and L.K.~Koopal,
``Determination of adsorption affinity distributions,''
\textit{J.\ Colloid Interface Sci.}\ \textbf{135}, 410--426 (1990).

\bibitem{Cernik1995}
M.~Cernik, M.~Borkovec, and J.C.~Westall,
``Regularized least-squares methods for the calculation of discrete
and continuous affinity distributions for heterogeneous sorbents,''
\textit{Environ.\ Sci.\ Technol.}\ \textbf{29}, 413--425 (1995).

\bibitem{Umpleby2001}
R.J.~Umpleby, S.C.~Baxter, Y.~Chen, R.N.~Shah, and K.D.~Shimizu,
``Characterization of molecularly imprinted polymers with the
Langmuir--Freundlich isotherm,''
\textit{Anal.\ Chem.}\ \textbf{73}, 4584--4591 (2001).

\bibitem{Umpleby2004}
R.J.~Umpleby, S.C.~Baxter, A.M.~Rampey, G.T.~Rushton, Y.~Chen, and
K.D.~Shimizu,
``Characterization of the heterogeneous binding site affinity
distributions in molecularly imprinted polymers,''
\textit{J.\ Chromatogr.\ B} \textbf{804}, 141--149 (2004).

\bibitem{Vieth1976}
W.R.~Vieth, J.M.~Howell, and J.H.~Hsieh,
``Dual sorption theory,''
\textit{J.\ Membr.\ Sci.}\ \textbf{1}, 177--220 (1976).

\bibitem{Paul1976}
D.R.~Paul and W.J.~Koros,
``Effect of partially immobilizing sorption on permeability and the
diffusion time lag,''
\textit{J.\ Polym.\ Sci.}\ \textbf{14}, 675--685 (1976).

\bibitem{Tikhonov1977}
A.N.~Tikhonov and V.Y.~Arsenin,
\textit{Solutions of Ill-Posed Problems},
Winston, Washington, DC (1977).

\bibitem{Hansen1998}
P.C.~Hansen,
\textit{Rank-Deficient and Discrete Ill-Posed Problems},
SIAM, Philadelphia (1998).

\bibitem{Hansen1992}
P.C.~Hansen,
``Analysis of discrete ill-posed problems by means of the L-curve,''
\textit{SIAM Rev.}\ \textbf{34}, 561--580 (1992).

\bibitem{Efron1979}
B.~Efron,
``Bootstrap methods: Another look at the jackknife,''
\textit{Ann.\ Statist.}\ \textbf{7}, 1--26 (1979).

\end{thebibliography}

\end{document}
